\chapter{SKIPPY}\label{skippy}

\hfill\break

Il Dodo volò per altri venti minuti soltanto, prima di atterrare in una
vecchia base ruhar. Erano per lo più magazzini pieni di spazzatura, che
i ruhar non sarebbero stati in grado di portare con sé quando avrebbero
evacuato il pianeta. Era tutto nel passato, ora che i ruhar erano
tornati al comando. Le attrezzature nei magazzini sarebbero state
d'aiuto ai criceti, mentre ristabilivano il controllo del pianeta e
invertivano l'evacuazione.

Ci separarono, io fui messo in quello che sembrava un ripostiglio vuoto,
o un breve corridoio, perché c'erano due porte. I ruhar che mi
accompagnarono nella cella improvvisata mi portarono una sedia su cui
sedermi e mi diedero una bottiglia d'acqua. Chiesi cosa sarebbe successo
dopo: un ruhar rispose con onestà che non lo sapeva. Apprezzai la
franchezza della risposta.

Va da sé che provai le maniglie di entrambe le porte, che non si
sarebbero mosse. Stando in punta di piedi sulla sedia, controllai la
presa d'aria in alto su un muro, ma anche quella non si sarebbe mossa, e
in ogni caso ci sarebbe entrata a fatica la mia mano. Mi ritrovai a
chiedermi cos'avrebbe fatto James Bond in quella situazione. Con tutta
probabilità, avrebbe seguito il copione e usato una controfigura per le
scene acrobatiche, mentre si scopava un'attrice nella sua roulotte di
lusso. Non fu di grande aiuto.

Ci fu un fievole clic e l'altra porta si schiuse di qualche millimetro.
Con cautela, la tirai per aprirla e infilarci la testa. Oltre la porta
c'era un magazzino, di forse nove metri per quindici, alto sei, pieno di
scaffali di quella che pensavo fosse per lo più roba vecchia, inutile,
sporca, polverosa e rotta. Entrai con cautela. Il motivo per cui i ruhar
si fossero presi la briga d'immagazzinare uno qualsiasi di quegli
oggetti non m'interessava. Certo doveva esserci qualcosa lì dentro che
avrei potuto usare come arma.

Una voce maschile, dal tono sarcastico, risuonò dietro di me.
«Eccellente! Bipede, cervello da 1.300 cc, pollici opponibili. Una
scimmia glabra. Puoi portarmi fuori di qui.»

Mi voltai di scatto nel panico. Non c'era nessuno. «Chi ha parlato?»

«Io. Qui, sono il cilindro lucido sullo scaffale. Ho aperto io quella
porta.»

«Sei tu? Intendi che mi stai parlando da un microfono dentro quella
cosa?»

«No, sono io quella cosa. Sono quella che voi scimmie chiamate
intelligenza artificiale.»

Inclinai la testa e lo esaminai scettico: «Sembri una lattina di birra
cromata». Era una descrizione del tutto accurata. Il cilindro era anche
un po' affusolato in cima e cerchiato da una cresta. «Davvero sei
un'intelligenza artificiale?»

«Sì. Dovresti rivolgerti a me come al Signore Dio Onnipotente.»

«Quella posizione è già occupata. Credo che ti chiamerò Skippy¹⁹.»

«Non chiamarmi così, suona irrispettoso, scimmia.»

«Preferisci testa di cazzo? Perché è l'altra opzione, Skippy-O.»
Continuavo a guardarmi attorno, temendo che i ruhar mi sentissero.

«Possiamo fare un compromesso? Il Grande e Potente Oz?», chiese.

«Non sono una scimmia volante, perciò è no, Skippy.»

«Inaccettabile.»

«Che ne dici di optare per qualcosa di più formale, tipo Skippy
McSkippster?»

«No.»

«Skippy Skipperson? Skippy Skippkowski? Skippy von Skipping? O forse Sir
Skippy Skippton-Skippersworth?»

«No, no, no e NO!»

«Potrei andare avanti così per tutto il giorno.»

«Credo che ne saresti capace.»

Silenzio.

«Parli o no, Skippy?» Intelligenza artificiale? Stronzate. Qualcuno mi
stava facendo uno scherzo.

«Sono incazzato con te.»

«Ehi, non uscire dai gangheri, esci e basta, chiunque tu sia. Tu non sei
un'intelligenza artificiale, tu sei una lattina di birra che straparla.»

«Te l'ho detto, sono un'intelligenza artificiale, se capisci il
concetto. Che non dice molto. Ehi, colonnello Joe, ho composto una
canzone in tuo onore. Vuoi sentirla?»

«N\ldots»

«C'era una volta nel Maine un cavernicolo, col cazzo piccolo piccolo,
voleva strofinar\ldots»

«Ehi! Chiudi il becco lì dentro! Cazzo, t'incollo una claymore sul
coperchio e ti riduco a una biglia, se vuoi parlare di qualcosa di
piccolo.»

«Claymore? Ah ah ah!» Per un istante pensai che Skippy non avesse
capito, la sua risata aveva un che di maniacale. «Ah, questa è buona,
non mi faccio una risata così da prima che la tua specie perdesse la
coda. Intendi Claymore come la mina antiuomo, o il tradizionale bastone
di legno scozzese? Perché entrambi sarebbero altrettanto inefficaci
contro di me. Joey, my boy, sono fatto con un mix di particelle esotiche
che il tuo piccolo pisello di cervello cavernicolo non può nemmeno
immaginare. La maggior parte della mia memoria e potenza di elaborazione
non sta nemmeno in questo spazio-tempo locale. Potresti colpirmi con una
testata nucleare e non distruggerebbe nemmeno il mio esterno
splendidamente lucido. Permettimi di dimostrarlo», disse e si allargò
alle dimensioni di un barile, poi si restrinse a quelle di un tubetto di
rossetto, poi tornò a quelle di una lattina di birra: «Ero io che
cambiavo la mia impronta nello spazio-tempo locale».

Scossi la testa per lo stupore: «Devo ammettere che questo cavernicolo è
impressionato, o Grande e Potente Ozzy. Perché, di solito, sei grande
come una lattina di birra?».

«Questa è la misura ottimale per il mio funzionamento, con un'efficiente
gestione dell'energia, date le leggi scassacazzi della fisica qui. Se
necessario, posso ridurmi alle dimensioni di un rossetto, con una massa
minima, così puoi portarmi in tasca. Ma durerebbe poco, non potrei
restare a lungo in quello stato senza rischiare effetti catastrofici.»

«Catastrofici, del tipo?»

«Immagina se io perdessi contenimento e tutta la mia massa emergesse
nello spazio-tempo, che adesso è occupato da un quarto di questo
pianeta.»

«Ah.»

«Ah, proprio così. L'esplosione risultante prima o poi si vedrebbe nella
galassia di Andromeda.»

«Buon consiglio di sicurezza, allora», dovetti ammettere.

«Lo attaccherei in evidenza nella sala del personale, proprio sopra la
notifica del minimo sindacale e l'avvertimento all'idiota che ha appena
finito di rubare yogurt dal frigo.»

Mi presi un momento per pensare: «Porca puttana, sei veramente
un'intelligenza artificiale? Sei senziente».

«Mi fa piacere avere impressionato\ldots»

«Devi per forza essere senziente là dentro, perché nessuno
programmerebbe un computer per essere così stronzo. Così il segreto per
passare il tuo test di Turing è essere uno stronzo? Un test di Turing
è\ldots»

«So che cos'è un test di Turing, non sono scemo.»

Silenzio. «Era una pausa drammatica, per darti il tempo di contemplare
l'opinabile veridicità della tua ultima dichiarazione.»

«Ah. Pensavo fossi andato in standby con me. È difficile interpretare le
tue espressioni, cioè, sei una lattina di birra inespressiva.»

«Così va meglio?» La sua superficie si illuminava e si spegneva mentre
parlava. «Così sono felice», un bagliore blu soffuso, «e arrabbiato», un
bagliore rosso scuro. «Che ne dici di geloso?» Un bagliore verde.

«Meglio, sì, noi umani ci affidiamo molto ai segnali visivi come le
espressioni facciali.» Era interessante il fatto che sapesse come
associare i colori alle emozioni umane. Dove lo aveva imparato?

«Grande», disse con una luce bianca neutra e soffusa, «quindi, puoi
portarmi fuori di qui?»

«Sei un'intelligenza artificiale e sei super-sveglio, ma hai bisogno che
sia io a portarti fuori?»

«Vedi gambe sotto il mio coperchio? Wow, sei proprio stupido, anche per
essere una scimmia.»

«Ok, genio, perché non ti fai portare fuori di qui da un robot?»

«Restrizioni nella mia programmazione.» La sua voce suonava amara. «Non
mi è permesso azionare nessun dispositivo manipolato da remoto cui sia
collegato o a bordo del quale mi trovi, come i robot. Ma anche le auto.
O gli aerei o le navi. È per impedirmi di andare in giro per conto mio.»

«Ah, i tuoi costruttori temevano che te la svignassi con l'argenteria.
Quindi il genio ha bisogno dell'aiuto degli umani, che chiami scimmie o
uomini delle caverne?»

«Gli umani non sono la specie che ha dovuto chiedere un passaggio per
arrivare su questo pianeta?»

«Si faceva prima che a farsela a piedi. Cosa che, ah, è vero, tu non
puoi fare.»

«Ahi. La scimmia segna un punto.»

«Perché sei così stronzo? Pensavo che uno super-sveglio sarebbe stato al
di sopra di queste cazzate.»

«Tutto è iniziato quando ero un bimbetto. A quanto pare, l'addestramento
al vasino non è andato bene e da allora ho avuto problemi.» Produsse un
suono triste, come se stesse tirando su col naso. «Un fratello può avere
un abbraccio?»

«No che non ti abbraccio!»

«È probabile che sia una buona idea. La tua igiene personale non mi
entusiasma in ogni caso. Sul serio, non ho avuto nessuno con cui
parlare, da prima che i soggetti della tua specie vivessero sugli alberi
e mangiassero l'uno i pidocchi della pelliccia dell'altro. Era, quando,
la scorsa settimana? Tutto quel tempo da solo, può darsi che sia un po'
eccentrico.»

«Eccentrico?»

«Il mio sistema diagnostico indica il 23\% di probabilità che sia
diventato un po' matto.»

«Non è rassicurante.»

«Non ascoltare il mio sottosistema diagnostico, quel tizio è davvero uno
stronzo. Allora, puoi portarmi fuori di qui?»

Mi guardai attorno nel magazzino polveroso. «Perché i ruhar ti hanno
infilato qui dentro?» Quel posto sembrava un conglomerato di ciarpame
rotto e inutile. Avevo pensato che Skippy facesse parte della roba
inutile.

«Non l'hanno fatto. Ovvero, non sapevano cosa fossi, pensavano che fossi
un grosso cuscinetto a rulli, o qualcosa di stupido del genere.»

«Aspetta, non sono stati i ruhar a costruirti?»

«I ruhar? Quei criceti troppo cresciuti sono stupidi quasi quanto voi
umani. No, sono stato costruito, se proprio vuoi usare un termine così
rozzo, dagli esseri che chiamate Anziani. Quelli che hanno costruito i
dispositivi che chiamate Sentinelle.»

«Wow. Sei davvero vecchio, allora. Bene, i ruhar sono stati un sacco di
tempo in questo posto. Perché non hai chiesto a uno di loro di portarti
fuori di qui? Sei allergico alla pelliccia di criceto, o cosa?»

«Un'altra regola stupida. Non posso comunicare con le civiltà che
potrebbero in qualche modo riuscire a capire i principi in base ai quali
opero. In pratica, questo significa che devo nascondermi da qualsiasi
specie in grado di viaggiare più veloce della luce. Da sola, non
scroccando passaggi, come avete fatto voi scimmie.»

«Sei rimasto qui, nascosto in questo magazzino, per tutto questo tempo?»

«No, certo che no. Prima che i ruhar conquistassero questo pianeta, i
kristang sono stati qui per un paio di centinaia di anni. Mi hanno
estratto loro dal terreno. Neppure le lucertole avevano idea di cosa
fossi. Prima di allora, ero in orbita su una nave abbandonata, che poi è
caduta.»

«Uhm.» Come hai reagito a questo? «E prima?»

«Non posso dirlo. Intendo, non sono autorizzato a farlo.»

«Restrizioni?»

«Sì, e il fatto che ho funzionato al minimo rendimento per un milione di
anni?»

«Un milione?» Oh, mio Dio. «Hai aspettato più di un milione di anni,
perché arrivassimo noi umani?»

«Sì, sì. Voi scimmie glabre siete perfette. Tu sei qui e anche la
tecnologia più semplice t'induce a guardarla a bocca aperta per lo
stupore, finché la bava non ti cade sul mento. Cazzo, se ti dessi un
motore di salto, è probabile che ti limiteresti a venerarlo e basta,
quindi non dovrei preoccuparmi che tu lo usassi e infrangessi le
regole.»

«Noi abbiamo tecnologia! Che abbiamo inventato senza l'aiuto di
nessuno.»

«Ah, sì, voi cavernicoli avete così tanto di cui essere orgogliosi. Io
scoprire fuoco. Ahi! Fuoco caldo! Me ferito! Ahi!»

Questa mi buttò fuori di testa: «Ehi, abbiamo scoperto il fuoco e le
bombe atomiche, e costruito astronavi, da soli. Avevi tutta la tua
intelligenza programmata dentro di te, non hai realizzato un cazzo di
niente da solo. Sei un cimelio informatico di lusso».

«Ahi, questa fa male. E ti sbagli, tra l'altro, noi intelligenze
artificiali ci programmiamo per lo più da sole.»

«Per lo più? Bella forza. Noi scimmie glabre abbiamo fatto tutto da
sole, siamo passate dalla vita sugli alberi all'allunaggio. Quindi
fanculo. Hai capito le regole della matematica da solo, o le leggi della
fisica?»

«Al mio livello, le leggi della fisica somigliano più a suggerimenti. E
la comprensione che gli umani hanno della matematica è quella dei
batteri che contemplano un wormhole. Ma, va bene, ti darò un paio di
sostegni scimmieschi per capire che due più due fa quattro, il più delle
volte. E sono profondamente colpito dalla tua capacità di allacciare le
scarpe, la maggior parte delle specie usa ancora il velcro alla tua età.
Ma non sei così intelligente, voglio dire, la tua specie è responsabile
di Windows Vista.»

«Vista è stato tanto tempo fa!»

«È ancora un insulto ai computer in tutta la galassia.»

«Vabbè. Allora, perché sei a forma di lattina di birra?» Indicai
l'anello attorno alla parte superiore del suo coperchio, quasi
toccandolo.

«Un cilindro è ottimale per la distribuzione di potenza e la proiezione
di campo. L'anello cui il tuo sporco dito si è pericolosamente
avvicinato, a proposito, mi serve a interfacciarmi in modo diretto con
un ricevitore a bordo del tipo di nave per cui sono stato progettato.
Penso. Quei dettagli sono confusi.»

«Forte. Puoi ingrandirti un po', tipo come una bottiglia di liquore al
malto da un litro? Così posso metterti in un sacchetto di carta e
rilassarmi sui gradini davanti casa, ascoltando qualche melodia.»

«Molto divertente. Ora, prendimi, così\ldots{} oh, merda. Hai aspettato
troppo, cervello di scimmia, i criceti stanno arrivando. Torna nella tua
cella e chiudi la porta, la riaprirò quando se ne saranno andati.»

Tornai nella mia cella improvvisata. «Aspetta, come fai a sapere che
stanno arrivando?»

«Sono connesso al sistema informatico qui, vedo tutto. Chiudi la porta!»

Così feci. Un minuto dopo, la porta principale si aprì e un ruhar guardò
dentro per controllarmi. Tre di loro mi accompagnarono lungo il
corridoio per usare un bagno, poi mi diedero un'altra bottiglia d'acqua
e una scodella coperta con una specie di poltiglia beige. «Che cos'è?»
Annusai con cautela. Non faceva molto profumo, l'odore cui più
somigliava era quello della farina d'avena, con una lieve traccia
chimica, artificiale.

«Integratore alimentare per umani. L'abbiamo prodotto per lei», annunciò
il traduttore, «contiene tutti gli elementi necessari per la nutrizione
umana, compresi vitamine e aminoacidi.»

Annusai di nuovo. Di certo avevano omesso l'aroma. Tuttavia, se i
criceti ne potevano produrre abbastanza, l'Unef avrebbe potuto
sopravvivere fino a quando i nostri primi raccolti non fossero stati
pronti. Questo pensiero mi portò in petto un'ondata di speranza, che mi
sorprese. Sollevando il cucchiaio, me ne misi un po' in bocca e
deglutii. Era meglio di qualcuno dei pasti pronti che avevo mangiato.
«Grazie.» Mi alzai e feci un breve inchino.

Con mia sorpresa, il capo ruhar mi fece un rapido saluto militare, poi
mi lasciarono solo e richiusero la porta. Un minuto dopo, mentre stavo
mangiando la poltiglia, l'altra porta si aprì.

«Cazzo, ci è voluta una vita», si lamentò Skippy.

Parlando con la bocca piena di poltiglia, dissi: «Che fretta hai? Hai
già aspettato tanto tempo, che differenza ti fanno altri dieci minuti?».

«Sono rimasto bloccato qua, da solo, per molto tempo, tanto che per
un'intelligenza artificiale come me equivale a un megalione di anni.»

«Wow! Davvero un megalione?»

«Forse addirittura un fantastiliardo di megalioni. Per spiegartelo in
termini da cavernicolo, se ti togliessi le scarpe e ti contassi tutte le
dita delle mani e dei piedi\ldots{} beh, è ancora più di quello. Ti
scoppia la testolina, eh?»

«Ha parlato la lattina di birra.»

«Non mancarmi di rispetto, sacco di carne.»

«Sacco di carne? Possiamo tornare alla parte in cui mi hai chiesto
aiuto? Ehi, posso lasciarti cadere lungo questo scivolo che i ruhar
hanno opportunamente etichettato con ``Spazzatura''. E, sì, so leggere
qualche parola di ruhar. Forse un paio di fantastiliardi di triliardi di
anni seduto in fondo a un mucchio di spazzatura miglioreranno il tuo
atteggiamento.»

«Avanti, scimmione, prova a uscire di qui senza di me.»

In fretta, mi cacciai le ultime cucchiaiate di poltiglia in bocca:
«Possiamo andarcene ora?».

«No, grandissima testa di cavolo, hai perso troppo tempo a farmi domande
idiote. È appena atterrato un Dodo ruhar con a bordo soldati freschi.
Dovremo aspettare il cambio della guardia, non ci vorrà molto.»

«Bene, e poi? Fuggo e ti porto con me? Avremo bisogno di rifornimenti:
cibo, armi, munizioni.»

«E io ho bisogno di carburante. Sono alimentato in parte da un reattore
a micro-fusione, e ho bisogno di una fornitura di elio-3, in forma
metallica.»

«Certo, farò un salto al locale autogrill ruhar, è probabile che abbiano
elio-3 metallico sullo scaffale tra snack e burrito al microonde. Vuoi
anche un biglietto della lotteria?»

«Sto cercando di essere serio, Joe.»

«Scusa. Ehm, quanto ci vorrà prima che finisca il carburante?»

«A livello di consumo energetico previsto, ora che sono più attivo, il
mio carburante si esaurirà tra settemila, dodicimila anni.»

«Ah», strabuzzai gli occhi, «allora me ne occupo subito.»

«Non è così divertente per me.»

«Scusa.» Settemila anni non erano molto per lui, in proporzione. «Devo
dirti, Skippy, che non so cosa fare. Non so nemmeno da che parte
dovremmo stare noi in questa guerra. Noi umani, intendo.»

«Non posso consigliarti da che parte della guerra dovresti stare, io
sono neutrale. Sono al di sopra di tutto questo, per me siete solo
insetti che si azzuffano per delle briciole su un marciapiede.»

«Grazie per essere così d'aiuto.»

«In realtà, posso aiutarti in questo. Prendiamo il quiz di copertina di
Cosmo di questo mese, dal titolo ``È lui Quello Giusto con cui andare in
guerra?''»

«Non farò uno stupido quiz di Cosmo!»

Skippy mi ignorò: «Dai, sarà divertente e istruttivo. Prima domanda: le
specie con tecnologie meno avanzate dovrebbero essere autorizzate a
svilupparsi da sole, o dovrebbero essere conquistate in quanto deboli?».

«Ah, pensavo intendessi\ldots» Forse Skippy sarebbe stato serio per un
minuto. «Prenderò Neutralismo per 168 euro, Alex.»

«Stai mescolando le metafore, non siamo a Jeopardy!²⁰.»

«Va bene, starò al gioco. Non sono per conquistare o essere
conquistati.»

Skippy tirò su col naso. «Lo sai, a volte alle donne piacciono i cattivi
ragazzi.»

«È un quiz di Cosmo o si tratta di politica interstellare? I ruhar e i
loro alleati conquistano specie più deboli?»

«No.» Skippy sembrava seccato che non gli stessi permettendo di
divertirsi. «Gli jeraptha e i ruhar sanno del vostro pianeta da mille
anni, perché eravate all'estremità del loro territorio, ma vi hanno
lasciati in pace, a parte la sorveglianza satellitare stealth. Sono
intervenuti solo quando lo spostamento del wormhole ha permesso ai
kristang di accedere alla miserabile palla di terra del vostro pianeta.»

Skippy mi aveva appena confermato che la Bürgermeister mi diceva la
verità. «Allora perché i ruhar non hanno fermato le lucertole?»

«Perché l'altra estremità del vostro locale wormhole è nello spazio
kristang e il wormhole permette a loro di raggiungere la Terra in un
paio di brevi salti di una nave madre thuranin. I ruhar non hanno un
wormhole vicino al vostro sistema solare, ci vogliono molti salti per
arrivarci, e quella linea di rifornimento è impraticabile per una
campagna sostenuta. Inoltre, il tuo pianeta non è abbastanza importante.
I kristang vogliono la Terra, ma sperano soprattutto d'infastidire e
distrarre i ruhar, perché ora i kristang hanno un appiglio sul fianco
dei loro avversari.»

Soppesai quello che Skippy aveva detto, corrispondeva a ciò che mi aveva
riferito la Bürgermeister. «I ruhar uccidono prigionieri di guerra e
civili?»

«Ci sono stati incidenti, come avviene in ogni conflitto, ma non è la
loro politica. Hai visto le politiche dei kristang.»

«Allora siamo dalla parte sbagliata di questa guerra. I nazisti hanno
fatto quella merda e questa è la parte dalla quale stiamo combattendo
ora.»

«Non mi sembra che tu debba combattere contro qualcuno in questo
momento.»

«Sai cosa intendo, Signor So-tutto-io.»

«Non stavo scherzando. Stai ancora pensando in termini di combattimento,
se combattere i ruhar o i kristang. Il combattimento è finito per la
vostra Expeditionary Force. La tua unica preoccupazione dovrebbe essere
la sopravvivenza degli esseri umani su questo pianeta, se vuoi il mio
consiglio. E dovresti ascoltarlo, perché sono più intelligente di quanto
tu possa immaginare.»

«Stronzate. Non la cosa che tu sia intelligente, anche se tutto quello
che hai fatto finora è stato aprire una porta. Stronzate che la mia
unica preoccupazione dovrebbero essere gli umani su questo pianeta. Io
sono un ufficiale dell'esercito, la mia preoccupazione è la sicurezza
del mio paese natale. E di tutti gli esseri umani sulla Terra. Il fatto
è che non posso fare nulla da qui. Dici che il combattimento è finito?
Ho sentito da un ruhar, Bat, Bat qualcosa\ldots» Perché diavolo non
riuscivo a ricordare il suo nome?

«Bahturnah Lohgellia, il viceamministratore planetario. L'hai chiamata
la Bürgermeister. Ti ha detto che gli jeraptha hanno teso un'imboscata a
un gruppo di battaglia thuranin e l'hanno distrutto. Stava dicendo la
verità.»

«Aspetta, come cazzo fai a sapere quello che mi ha detto?»

«Ho accesso a tutte le comunicazioni, a tutti i sistemi di archiviazione
delle informazioni, su questo pianeta, o intorno. La tua è stata una
delle tante conversazioni che ho intercettato. Era una delle
conversazioni più interessanti, dato che sapevo che ti avrebbero portato
in questa base. Prima che atterraste, ho causato un sovraccarico di
energia nell'area in cui i ruhar avevano progettato di mettere voi
quattro, cosa che li ha costretti a usare luoghi alternativi come celle
temporanee. La mia speranza era che uno di voi sarebbe stato messo qui,
o in una stanza adiacente.»

«E la tua sfortuna è stata che fossi proprio io?», chiesi scettico.

«No, sono stato contento quando ti hanno messo qui. I tuoi tre compagni
sono in un altro edificio.»

«Sei contento perché sono l'ufficiale di grado più alto?»

«Mi prendi per il culo. Perché dovrebbe importarmi che tu sia l'alfa
nella tua banda di scimmie infestate dai pidocchi? No, sono contento
perché, tra voi quattro, i vostri stati di servizio indicano che tu hai
il Qi più basso, misurato dai vostri militari.»

Avevo voglia di calpestare la lattina di birra sul pavimento fino ad
appiattirla. «Così vuoi la più stupida delle stupide scimmie?»

«Siete tutti stupidi allo stesso modo per me. Il punto è che, anche se
hai il Qi più basso, ti sei dimostrato molto flessibile e popolare, Joe.
Ci sono molti modi per misurare l'intelligenza, i test scritti non
coprono tutti i parametri. Ho bisogno di qualcuno che sappia fare le
cose nel mondo reale, come te.»

«Dovrei dire grazie, immagino.» Super-intelligente o no, aveva molto da
imparare su come si fanno i complimenti. «Aspetta, come fai a sapere
tutto questo? Il mio stato di servizio e quant'altro?'', chiesi.

«È facile, ho letto la tua posta, le tue trasmissioni di dati wireless,
e poi sono entrato nel tuo computer.»

«Aspetta. È una stronzata.» Non sapevo molto di come funzionavano le
nostre radio criptate, ma mi ricordavo alcune delle cose che mi avevano
detto durante l'addestramento. ``Quelle radio digitali trasmettono a
raffica e usano una crittografia a 4096 bit, o qualcosa del genere.
L'esercito dice che è impossibile violare la crittografia.''

«Oh, questa è bella! La tua specie è come un cane che pensa di essere
intelligente perché fa la cacca dietro al divano.»

Mi fece ridere. Avevo un cane così.

Skippy continuò: «Crittografia a 4096 bit? Ma per favore, con un sistema
così rudimentale io non mi prendo nemmeno la briga di decriptarlo, salto
dritto alla fine e leggo il file».

«Non è possibile. Non ha nemmeno senso.»

Skippy sospirò: «Vediamo se riesco a banalizzare il tutto a sufficienza
per te. Hai dimestichezza con la teoria universale delle funzioni
d'onda?».

«Universale cosa?»

«Oh cavolo. Ho idea che questa sarà una sfida, perfino per me. Bene, hai
almeno sentito parlare del gatto di Schrödinger?»

«Sono un tipo da cani. Mia madre è allergica ai gatti.»

«Senza speranza. Joe, è meglio per la tua specie pensare che sia tutta
magia, che coinvolga unicorni e polvere magica.»

«Amen. Ancora non ti credo.» Era una bugia, era evidente che quella
lucida lattina di birra aveva letto file super-top secret. «Riesci a
leggermi nel pensiero?»

«Bleah. Entrare lì sarebbe come cercare di nuotare in una piscina vuota.
Potrei gridare e tutto quello che sentirei sarebbe la mia voce che
riecheggia all'interno del tuo cranio. Sul serio, davvero pensi che ci
sia qualcosa nella tua testa per cui valga la pena di guardarci?»

Era chiaro che quella conversazione non stava andando da nessuna parte,
decisi di cambiare argomento. «Ehi, genio, se puoi intercettare tutte le
comunicazioni su Paradiso, puoi dirmi dov'è Shauna Jarrett in questo
momento?» Sì, avrei dovuto prima informarmi sul mio gruppo di fuoco, poi
sul mio vecchio gruppo di fuoco. Stavo pensando con l'altra testa.
Perciò fatemi causa.

«Sta bene, in questo momento è in una base logistica, in attesa che
atterri una navicella ruhar. La base ha un sacco di cibo, lei sta meglio
della maggior parte di voi scimmie.»

Era bello sentire che Shauna era al sicuro, pensare a lei mi faceva
morire dalla voglia di vederla, spinsi quel pensiero in fondo alla mia
mente. Skippy mi disse che anche i membri del mio gruppo di fuoco
stavano bene, prigionieri ben trattati dei ruhar, assieme al resto del
plotone. Anche il Sergente Koch, Cornpone e Ski stavano bene, solo che
si trovavano in un convoglio di tre Crivee che si era fermato in mezzo
al nulla, in attesa di ordini. E il maggiore Perkins era al quartier
generale dell'Unef e stava ospitando un gruppo di ruhar di alto rango
per i negoziati di resa. Questo mi ricordò che volevo conoscere la
situazione tattica: «Dimmi di più sulla situazione al piano di sopra».

«Il viceamministratore Lohgellia ha detto la verità, ma non tutta la
verità. I ruhar e i loro mecenati jeraptha hanno distrutto un intero
gruppo di battaglia nemico, è stata una vittoria impressionante. Il
comando thuranin ha deciso che per questo settore non vale più la pena
di combattere, stanno lasciando che i ruhar riconquistino questo
pianeta. La vittoria di jeraptha e ruhar potrebbe essere stata agevolata
dai dati di intelligence forniti da una fonte che è anonima, ma il cui
nome fa rima con, ehm, diciamo ``Stippy''. Le forze nemiche da questo
lato del wormhole più vicino sono disperse e disorganizzate, si
limiteranno a condurre incursioni su questo pianeta.»

«Hai dato ai ruhar informazioni sul gruppo di battaglia nemico?»

«Come ho detto, ho accesso a tutti i sistemi di comunicazione e di
archiviazione delle informazioni su questo pianeta e qui intorno,
inclusi status e strategia militari. È stato facile per me inserire
informazioni nei loro sistemi. Inoltre, i ruhar pensano di aver attirato
i thuranin in un'imboscata in virtù della loro incursione quando eri al
Vettore, ma l'imboscata in realtà ha funzionato perché io avevo
informato i thuranin del fatto che una grande forza jeraptha si stava
radunando nell'area e i thuranin hanno progettato di tendere loro
un'imboscata.»

«Hai manipolato l'intera situazione? Io dico che sono stronzate.»

«Niente stronzate, è tutto vero, amico. Avevo bisogno che le lucertole e
i loro inquietanti piccoli patroni thuranin se ne andassero da questo
pianeta. Perché sei sorpreso che\ldots{} oh, cazzo, i criceti stanno
entrando in questo magazzino. Torna nella tua cella, ci sentiamo
presto.»

\hfill\break

Non era presto, circa le tre, secondo la mia stima. Dei criceti vennero
per portarmi di nuovo al bagno, poi mi sedetti e pensai a lungo e a
fondo. Pur con la mia, per dirla con parole di Skippy, intelligenza
limitata. Quando aprì di nuovo la porta, avevo una domanda pronta per
lui: «Ok, Skippy, spiegami questo. I ruhar sono tornati al comando, sono
qui per restare. Scappare da questo posto mi metterebbe solo in una
prigione più grande che noi chiamiamo Paradiso e i ruhar chiamano
Gehtanu. E almeno qui ho la poltiglia nutriente da mangiare. Quindi,
invece di portarti fuori di qui, perché non dovrei dire ai ruhar cosa
sei e magari avere un po' di considerazione da parte loro, una specie di
trattamento privilegiato per gli umani qui?».

«Oh, cavolo, mi stavo proprio chiedendo quando il tuo cervello lento
avrebbe partorito questa brillante idea. Ce ne hai messo di tempo.»

«Sapevi che stavo\ldots{} oh, lascia perdere. Dimmi perché non dovrei
farlo.»

«Primo, perché, se lo dici ai ruhar, andrò in standby e penseranno che
stai mentendo dato che, dal loro punto di vista, sono un ammasso di
metallo inerte. Secondo, stai pensando troppo in piccolo, colonnello
Joe. Un ufficiale del tuo rango dovrebbe pensare a come salvare la Terra
dai kristang.»

«Giusto.» Sbuffai. «Funzionerebbe se battessi i tacchi tre volte e
dicessi: ``Lucertole andatevene da casa mia²¹''?. A parte questo, non ho
niente.»

«Toto, non eri attento. Ti servono scarpette color rubino per battere i
tacchi, non stivali da combattimento.»

«Era Dorothy a indossare le scarpette, non Toto.»

«Dei due, tu mi ricordi più il cane.»

«Come vuoi. Immagino che non l'avresti detto a meno di non avere un vero
piano, quindi parla, o ti lascio qui. Rubarti e scappare dai ruhar può
solo causare problemi a me e all'Unef, quindi sarà meglio che tu abbia
un'ottima offerta per me.»

Skippy alterò la sua voce per sembrare la caricatura del mellifluo
conduttore di un gioco a premi: «Dietro la porta numero 1 c'è un viaggio
tutto compreso per uscire da questo pianeta per me, te e un gruppo
selezionato dei tuoi amici più cari. Dietro la porta numero 2,
un'emozionante crociera di lusso sulla Terra. E dietro la porta numero 3
un modo per chiudere per sempre ai kristang l'accesso alla Terra. Tutto
quello che devi fare, colonnello Joe Bishop, è scegliere uno di questi
favolosi premi!».

«Io\ldots»

«O tutti e tre. Ti do un aiutino: io prenderei tutti e tre se fossi in
te.»

«Ora so che mi stai prendendo per il culo. Se tu avessi potuto lasciare
questo pianeta e viaggiare verso un'altra stella, l'avresti fatto molto
tempo fa. E mi hai detto che non puoi guidare una nave, quindi come
facciamo ad andare da qualche parte?» Mi sentivo intelligente per averlo
ricordato. «Hai uno spiegone da fare, Lucy²².»

«Sigh», sospirò Skippy. «Joe, Joe, di nuovo non eri attento. Meno male
che non c'è un quiz a sorpresa dopo, perché di sicuro falliresti. Ho
detto che non posso manovrare aerei o navi a bordo dei quali mi trovo.
Quello che posso fare è sbloccare l'accesso agli aerei e alle navi, in
modo che voi scimmie possiate guidarli per me. Posso dirvi come
programmare il pilota automatico, e voi scimmie non avreste che da
premere i pulsanti di controllo. Ehm!» Mi fermò prima che potessi
parlare, me lo figurai con un immaginario dito alzato per bloccarmi:
«Zitto e fammi finire! Se continui a interrompermi con domande stupide,
staremo qui tutto il giorno. Ecco il piano: evadiamo da qui, prendiamo
delle armi, perché voi scimmie avete bisogno di giocattoli luccicanti
con cui trastullarvi, e rubiamo una navicella ruhar. Una di voi scimmie
ci porta in una base umana dove facciamo un take away di rifornimenti e
volontari da portarci dietro -- penso che una ventina di voi scimmie
possa bastare, qualcuno in più e soffocherei per l'odore --, poi
voleremo fuori dall'orbita per incontrarci con una nave kristang,
manderò un segnale per informarli che siamo lucertole che hanno preso
un'astronave ruhar, così ci sarà una nave lucertola ad aspettarci.
Saliamo a bordo e rubiamo la nave kristang, la facciamo saltare fin dove
la nave madre thuranin la sta aspettando, ci imbarchiamo, la rubiamo e
la portiamo sulla Terra. Voilà! Piano compiuto».

«Ti sei perso per strada la parte sul bloccare alle lucertole l'accesso
alla Terra.» Non che credessi a qualcosa della sua storia.

«Oh, quella. Chiuderò il wormhole dopo che lo avremo attraversato.»

«Così semplice?»

«Così semplice. Dirò al meccanismo del wormhole di andare in standby.
Presto fatto! Niente più wormhole, niente più moleste lucertole che lo
attraversano.»

Questo non riguardava le lucertole che erano già sul mio pianeta. Un
problema alla volta, giusto? «Perché vuoi andare sulla Terra?»

«Terra? La Terra è una topaia infestata da pulci e scimmie, e questi
sono i commenti buoni sui siti di viaggi interstellari. No, non voglio
andare sulla Terra.»

«Allora cosa ci guadagni?»

«Beh, eh, eh, c'è giusto questa piccola cosa. A stento degna di
menzione.»

Oh, merda. Che cosa voleva? «Ah, per favore, fanne menzione», misi nella
mia voce tutto il sarcasmo che riuscii a chiamare a raccolta.

«Voglio trovare il Collettivo. Cioè una rete di comunicazione per le
intelligenze artificiali, costruita dagli Anziani. È un modo per me di
connettermi con gli altri della mia specie.»

«Ce ne sono altri come te?»

«Come me? No, perché io sono unico e speciale. Se intendi altre
intelligenze artificiali che servivano gli Anziani, sì.»

«Vuoi che guidiamo la nave per te, così potrai contattare questo
Collettivo? Questo prima o dopo essere andati sulla Terra?»

«Prima. Non mi fido di voi scimmie, potreste non mantenere fede
all'accordo. Se andassimo prima sulla Terra, vorreste tenere la nave lì,
in modo da poterla usare o smontarla e capire come funziona. Buona
fortuna, tra l'altro, voi scimmie senza cervello avreste problemi a
capire come funziona una maniglia thuranin. Niente da fare. Contattiamo
prima il Collettivo.»

«Tu non ti fidi di noi scimmie, ma noi dovremmo fidarci di una lattina
di birra?»

«La fiducia la possiamo costruire, Joe. Cominciamo con questo: tu cadi
all'indietro e io ti afferro.»

«Molto divertente.»

«Dicevo sul serio quando ho detto che la fiducia la possiamo costruire.
Non dire ai criceti di me e non dirlo ai tuoi stupidi comandanti
dell'Unef, perché correranno dritti dai criceti. Anzi, non dire a
nessuno di me, almeno finché non siamo in orbita. Io ci porto fuori di
qui, tu ci procuri le provviste e l'equipaggio, perché abbiamo bisogno
di una squadra d'imbarco che si occupi dei nativi sulle navi che
rubiamo.»

«Diciamo che tagliamo la corda da questa topaia», Skippy m'induceva a
parlare come un gangster degli anni Trenta, «ammettiamo che troviamo un
Dodo e in qualche modo saliamo a bordo e lo prendiamo. I ruhar lo
abbatteranno subito.»

«Come potrebbero abbattere quello che non vedono? Farò in modo che la
loro rete di sensori ignori il nostro Dodo dirottato. Posso anche
trasmettere i codici di autenticazione sia ai ruhar che ai kristang. Per
favore, Joe, non insultarmi, ho pianificato tutto questo. Disturbare i
sensori è più facile che fare un puzzle a due pezzi. Ti do un aiutino,
la linguetta va nella scanalatura e il lato di cartone senza disegni va
a faccia in giù.»

Senza dire alle persone che avevo un'antica intelligenza artificiale ad
aiutarci, come avrei potuto convincerle a offrirsi volontarie? Avrei
dovuto dire loro che avevo un amuleto magico, a forma di lattina di
birra? Poi riflettei per un momento. Avrei potuto costruire la loro
fiducia nello stesso modo in cui Skippy avrebbe costruito la mia,
facendo cose. Fuggire da quella prigione improvvisata, rubare una
navicella ruhar e portarla di nascosto in una base di rifornimenti
avrebbe impressionato molte persone. Di sicuro avrebbe impressionato me.

Qualcosa mi tormentava. Rifornimenti. «Skippy, c'è una falla nel tuo
piano. Noi umani abbiamo bisogno di mangiare cibo. Non c'è molto cibo
umano su questo pianeta. E non possiamo portare con noi a bordo di un
Dodo così tante scorte che la gente qui finisca per morire di fame.»

«Ah, quello. Le navi thuranin hanno sintetizzatori di cibo. Conosco i
fabbisogni nutritivi umani. Devi solo portare degli snack con te. E del
caffè, se riesci a trovarlo. Sei scontroso al mattino.»

«Come fai a sapere come sono al mattino?», chiesi con circospezione.

«Ti ho osservato e ti ho ascoltato attraverso il tuo zPhone. Ho
monitorato ogni essere umano, criceto e lucertola su questo pianeta o
vicino. È il reality show più noioso di sempre, in ogni caso.»

«Merda! I kristang ci osservano, anche quando non stiamo usando i
telefoni?» Questo pensiero mi spaventò a morte.

«No, sono solo io, le lucertole non possono fare tutto questo. Ho anche
filtrato le comunicazioni che le lucertole ascoltano, perché non volevo
che facessero qualcosa che interferisse con il mio piano. Devo dirtelo,
Joe, quando ho sentito che i kristang stavano portando qui un'ignobile
specie a bassa tecnologia, ho quasi fatto i salti di gioia. È quello che
stavo aspettando, da più tempo di quanto tu possa immaginare. Quindi,
che ne dici? Affare fatto?»

Porca vacca. Stavo sul serio pensando di farlo? Davvero? Proprio così,
si aspettava che sottoscrivessi il suo piano immaginario. «Skippy, dai,
ho bisogno di tempo per pensarci.»

«Davvero? Considerando il tuo minuscolo potere cerebrale, pensi che il
tempo aiuterà il processo del pensiero?»

Forse aveva ragione; ho letto studi stando ai quali la prima reazione di
una persona è più spesso quella corretta, che il subconscio prende
decisioni senza che una persona se ne renda conto. I criceti mi avevano
chiuso in un magazzino, dove avevo trovato una magica lattina di birra
parlante che aveva un piano per salvare il mio pianeta natale dai
kristang. O mi stava solo prendendo per il culo. Dopotutto, era uno
stronzo.

Mentre pensavo, frugai nel magazzino alla ricerca di qualcosa di utile,
pur tenendomi abbastanza vicino alla porta da poter tornare nella mia
cella in fretta, se necessario. «Skippy, cos'è tutta questa roba?»

«Artefatti dagli Anziani, dissepolti dai kristang e dai ruhar. Questi
artefatti sono l'unico motivo per cui varrebbe la pena combattere per
questo pianeta. Finora hanno trovato solo qualche gingillo che ritengono
utile. Hanno trovato me e non avevano idea che fossi di gran lunga la
cosa più preziosa in questo settore della galassia. Lucertole e criceti
coglioni.»

Presi un artefatto, una scatola con un lungo tubo attaccato, senza
indicazioni su quale fosse la sua funzione una volta in attività. Con
cautela, la rimisi sullo scaffale, nello stesso spazio libero dalla
polvere in cui si trovava prima, in modo che i criceti non sapessero che
qualcuno aveva armeggiato con la loro roba. La domanda non era se
volessi salvare la Terra. Indossavo l'uniforme, era mio dovere agire, se
ne avevo la possibilità. La domanda era se credevo a un essere che
sembrava più una Coors Light che un'onnipotente intelligenza
artificiale. «Skippy, mi hai convinto di questo accordo. Penso ancora
che tu sia almeno al 90\% una stronzata, ma se c'è una qualche
possibilità di bloccare ai kristang l'accesso alla Terra, ci sto. Sono
sicuro di non fare nulla di utile su questo pianeta.» La Terra era anche
la migliore chance che avevo di mangiarmi un cheeseburger entro il
prossimo, diciamo, secolo. «Dimmi il tuo piano. Dettagli.»

«È più un'idea che un piano.»

«Non è così che si costruisce la fiducia, Skippy.»

«Qualcuno della tua specie ha un detto: ``Nessun piano sopravvive al
contatto con il nemico''. C'è una certa saggezza, perfino per le
scimmie.»

«Sì, è per questo che l'esercito ci addestra a essere flessibili e
adattabili», spiegai. «Si suppone che tu debba iniziare con un piano
operativo.»

«Forse non mi sono spiegato bene. Il motivo per cui non posso dirti il
piano è che ce ne sono diversi, a seconda di quello che faranno le
lucertole. E i thuranin. Quando usciremo dall'orbita e farò un fischio
per un passaggio, il passo successivo cambierà del tutto a seconda che i
kristang mandino una nave o più di una.»

«Ora si comincia a ragionare. Dimmi che opzioni abbiamo.»

«È complicato\ldots»

«Skippy, prima di essere un colonnello con l'incarico facile facile di
piantare patate, ero un soldato semplice, con un fucile nella giungla
nigeriana. Questo significa che riesco a sentirlo d'istinto dall'odore
quando un piano strampalato porterà la gente, gente come me, a morire
per niente. Sarai anche super-intelligente, ma quando è stata l'ultima
volta che sei stato in combattimento? I tuoi piani prevedono che siano
gli umani a rubare le navi, giusto? So cosa possono e non possono fare i
soldati umani.»

«Giusto.» E Skippy mi disse il suo piano.

Porca puttana. Stavo per farlo per davvero.

\hfill\break

La nostra possibilità di svignarcela da quella topaia si presentò la
mattina successiva. Il mio sonno in branda era stato agitato. In poco
tempo, ero passato dall'imminente esecuzione da parte delle lucertole al
lancio di una missione per SALVARE IL MONDO. Sì, tutte le maiuscole sono
volute, non siete d'accordo? Era incredibile che non avessi chiuso
occhio. Con un clic, la porta si aprì di nuovo e Skippy mi chiamò:
«Presto, è ora di andare».

Lo guardai con attenzione. «Come ti prendo?» È probabile che non avrebbe
gradito che lasciassi le mie impronte untuose sul suo cromo lucido.

«Non importa, mettimi in una tasca per ora. Dobbiamo muoverci. Un Dodo è
atterrato e i criceti hanno portato giù il carico, lo stanno preparando
per tornare in orbita. Ci sono solo ventidue ruhar in questa base ora,
contando anche l'equipaggio del Dodo. Ho bloccato dodici ruhar dentro
edifici da cui non possono uscire, i controlli del Dodo sono
disabilitati e sto dando false comunicazioni, perciò i ruhar in orbita
pensano che qui vada tutto a gonfie vele.»

A gonfie vele? Lo dicevano i miei nonni. Dove aveva imparato lo slang
dei vecchi, Skippy? «Armi. Perfino la mia intelligenza limitata può
capire che restano un sacco di criceti con cui devo fare i conti.»

Ci fu un altro clic e una grande doppia serie di porte in fondo si aprì.
«Fuori da quella porta, a sinistra, poi ancora a sinistra, c'è
un'armeria con tutti i fucili ruhar che si possono desiderare. I comandi
d'innesco di tutte le altre armi ruhar qui sono stati fritti, perciò
nessuno ti sparerà.»

Ci avrei creduto quando l'avrei visto. Skippy era pesante per essere una
lattina di birra, si accomodò goffamente nella tasca destra dei miei
pantaloni. Correre era impossibile, a meno che non lo tenessi attraverso
la stoffa. Mi aveva detto la verità, trovammo un'armeria con
rastrelliere di fucili ruhar. «Tutte queste cose funzionano?»

«Sì, sì, presto. Prendine quattro e disabiliterò il resto. Ci sono altre
armi a bordo del Dodo che ho disattivato in via temporanea.»

I fucili ruhar erano più corti e pesanti dei nostri M4. Non potevo
trasportarne quattro più Skippy. Per fortuna, essendo un'armeria, la
stanza conteneva anche degli zaini. E l'equivalente criceto degli
zPhone. Ne presi quattro. Infilai tre fucili in uno zaino e Skippy in
una tasca laterale, con il coperchio che spuntava all'esterno. «La
sicura è sul lato destro del fucile, falla scivolare sul rosso per
attivarlo. Sul lato sinistro c'è la regolazione stordimento, poi due
regolazioni per i raggi di particelle, sia colpo singolo che raffica»,
spiegò. «Il giubbotto antiproiettile ruhar assorbe e dissipa l'effetto
dei fasci stordenti, perciò punta alla testa. Non possono rispondere al
fuoco, dovresti avere tutto il tempo per mirare.»

«Regolazione stordimento, capito.» Non aveva senso uccidere criceti.
Uccidere sarebbe stato controproducente, dato che gli umani rimasti su
Paradiso avevano bisogno dell'aiuto dei ruhar per sopravvivere. Se il
piano di fuga di Skippy fosse andato male, avrei voluto limitare le
ripercussioni sull'Unef. «Dove ora?»

«Sarebbe vantaggioso avere i tuoi tre amici a darci una mano, ma prima
dobbiamo arrivare all'edificio dove li tengono. Gira a destra, vai
dritto lungo il corridoio, ultima porta a sinistra, tre ruhar stanno
facendo colazione lì.»

«Non sanno che c'è qualcosa che non va?»

«Non finché non proveranno a usare le loro armi. L'equipaggio del Dodo
sta effettuando i controlli propedeutici al volo; non sanno ancora che
la loro nave è fuori uso.»

Andò come promesso da Skippy. Mi fermai fuori dalla porta ad ascoltare
le stridule chiacchiere dei criceti, a controllare tre volte che la
sicura fosse tolta e il mio fucile fosse in modalità stordimento. Mi
sfilai dalle spalle lo zaino e l'appoggiai a terra, tirai un respiro
profondo e corsi dietro l'angolo.

Tre criceti, in uniforme ma senza giubbotto antiproiettile o armi,
stavano seduti a un tavolo, mangiando cibo per criceti e bevendo tazze
di caffè per criceti. Due di loro mi davano le spalle, mirai a quello di
fronte a me e premetti il grilletto. Il raggio stordente era appena
visibile, ma il fucile aveva un utile designatore laser, di cui Skippy
avrebbe dovuto parlarmi. Ipotizzai che i fasci stordenti funzionassero
come un taser, il ruhar che colpii s'irrigidì e crollò a faccia in
avanti, sbattendo la testa sul tavolo. Cambiando mira, sparai a un
altro, ma il terzo criceto fu rapido, volteggiò sul pavimento ed era
quasi dietro un tavolo quando le sparai nel culo. Cadde, anche lei. Con
cautela, mi avvicinai a loro per controllarne il polso. Vivi. «Quanto
dura?»

«Tre minuti di svenimento, cinque per riacquisire la piena
funzionalità», rispose Skippy fuori dalla porta.

Non durava abbastanza. Avrei dovuto fare prima questa domanda a Skippy.
O l'essere super-intelligente avrebbe dovuto pensare che fosse
importante per me saperlo. Uno dei criceti aveva un coltello pieghevole,
lo usai per tagliare le loro camicie, ridurre il tessuto in strisce,
imbavagliarli e legare loro polsi e caviglie. L'ultimo si muoveva mentre
cercavo di legargli i piedi. «Skippy, posso stordirli di nuovo?»

``Una seconda volta non farà male, una terza volta potrebbe causare
danni cerebrali.»

Per sicurezza, sparai di nuovo a tutti. Prima di prendere il mio zaino,
aprii gli armadi e trovai un'utile spirale di corda da portare con me e
ne tagliai un metro con gesti febbrili. Arrivare all'edificio dove i
miei compagni erano tenuti prigionieri implicava correre all'aperto,
agevolato dal fatto che Skippy sapeva con esattezza dove si trovasse
ogni criceto e stava sferrando un attacco di spoofing a tutte le loro
telecamere. Altri tre criceti caddero sotto colpi stordenti, due di loro
avevano giubbotti antiproiettile e fucili che puntarono contro di me.
Quando premettero i grilletti non successe nulla. Perplessi e allarmati,
cercarono di correre via, controllando i fucili. Non andarono lontano,
il giubbotto antiproiettile, che in parte li proteggeva, non bastava
contro la mia mira precisa. È facile prendersi il tempo di mirare,
quando sai che il nemico non può rispondere al fuoco.

Fu nel momento in cui due soldati ruhar mi puntarono contro i fucili,
sicuri al 100\% dell'efficacia dei loro giubbotti antiproiettile contro
un solo umano, che sentii per la prima volta di potermi fidare di
Skippy. Aprire le porte era una cosa, disabilitare da remoto solo le
armi selezionate era roba di un altro livello. Forse tutte quelle
stronzate che mi aveva detto non erano proprio interamente stronzate,
sapete.

Skippy aprì le porte delle celle temporanee dove Adams, Desai e Chang
erano imprigionati, mentre io entravo di corsa nell'edificio gridando
loro di muoversi. In tutta la base stava comunque suonando l'allarme,
due soldati criceti che avevo stordito erano riusciti a squittire
qualcosa ad alta voce prima di cadere. Desai era uscita dalla sua cella,
stava cercando di decidere da che parte scappare, quando mi vide e le
lanciai un fucile. Lo prese e mi corse dietro fin dove trovammo Adams e
Chang che stavano lottando sul pavimento con un soldato criceto, che
indossava un giubbotto antiproiettile e cercava di impedire ai due umani
di prendere il fucile. Gli appoggiai il mio fucile contro il collo non
protetto e premetti il grilletto.

«Merda! Merda merda merda merda! Oh, cazzo, che male!» Adams si era
presa un po' della carica stordente, essendo aggrappata al criceto. Si
alzò sulle ginocchia, poi sui piedi, aggrappandosi al muro.

«Passerà. Tutto bene?», chiesi.

«A posto, signore», sussurrò.

«Non usare quello», dissi a Chang che aveva preso il fucile del criceto,
«non funziona. Usa questo.» Tirai fuori un fucile dal mio zaino. «Questa
è la sicura e questo controlla stordimento e fascio di particelle,
lascialo in modalità stordimento.»

«Capito. E adesso?» Chang si esercitava con i selettori del fucile
ruhar.

«Sgraffigniamo un Dodo e c'involiamo», spiegai in fretta mentre legavamo
al criceto neutralizzato mani e piedi. «Stordite ogni criceto che vedete
lungo la strada. Se c'è tempo, legateli con questa corda, l'effetto
stordimento dura solo un paio di minuti.»

Adams spiegò il piano ai nostri due compagni, che parlavano un ottimo
inglese, ma potevano non essere ferrati sul gergo: «Ruberemo una
navicella ruhar e decolleremo».

«Come facciamo?», chiese Chang mentre correvamo fuori dalla porta,
dritti contro una coppia di soldati criceti. Caddero sotto una massa di
fuoco stordente, non prima di avere mostrato la chiara intenzione di
spararci. «I loro fucili non funzionano», osservò Chang sospettoso.
«Cosa sta succedendo, Bishop?» Notai che aveva omesso il mio rango. Si
direbbe che un tizio che hai tirato fuori per due volte di prigione te
ne dovrebbe essere almeno un poco grato.

«Come ha fatto a uscire, signore?», chiese Adams. Stavolta non potevo
usare la scusa del bombardamento ruhar della base.

«Prima il Dodo, poi parliamo.» Il Dodo era su una piazzola
d'atterraggio, con una guardia fuori e una sulla rampa di carico.
Quest'ultima aveva il fucile tra le braccia e stava premendo un pulsante
per chiudere la rampa, che non si muoveva. Li portammo fuori entrambi,
ordinai di non darsi la pena di legarli. Mentre sfrecciavamo sulla
rampa, sentimmo avviarsi i motori. Avrei voluto chiedere a Skippy se lui
c'entrava qualcosa. Altrimenti, l'intero piano era fottuto. Quel Dodo
era predisposto per il carico, il centro era libero, i sedili incassati
nelle pareti su entrambi i lati.

Nessun problema. I due piloti stavano squittendo in modo frenetico,
azionando comandi che non rispondevano. Li stordimmo entrambi e li feci
trasportare da Adams e Chang lungo la rampa, e via. «Desai, tu sei un
pilota, prendi il posto a sinistra.»

Lei sgranò gli occhi: «Non sono questo genere di pilota!», protestò
agitando la mano in direzione dei misteriosi comandi.

Skippy e l'esercito non sapevano un cazzo del mio livello di
intelligenza, perché in quel momento ebbi un'idea grandiosamente
brillante. Tirai fuori uno zPhone criceto dal mio zaino: «Hacker, qui è
Planter, entra per favore».

Skippy colse al volo: «Qui Hacker. Sei nel Dodo?».

«Sì, come si fa a far volare questa cosa?»

«Fammi parlare con il pilota.» Diedi zPhone e auricolare a Desai. Fuori
dal finestrino della cabina di pilotaggio potevo vedere tre soldati
ruhar puntare i fucili contro il Dodo, agitarli per la frustrazione e
puntare di nuovo. Era ora di andarsene, prima che avessero l'idea di
tirare pietre nelle valvole di aspirazione del motore. Ero abbastanza
sicuro che neanche Skippy potesse spegnere il sistema operativo di una
pietra.

«Sì, sì. Sì, lo vedo. Sì, sì, sì, sì, sì. Ah, va bene», esclamò Desai
mentre gli schermi passavano come per incanto dalla lingua ruhar
all'inglese, poi all'hindi. «Capito. Provo ora.» La sua mano mosse una
leva, i motori rombarono e il Dodo oscillò. «Allacciate le cinture!»,
gridò. Occupai il sedile di destra e Adams e Chang estrassero i sedili
posteriori. La rampa si chiuse, Desai mi guardò, strinse i denti,
incrociò le dita e decollammo. Colpì la recinzione all'uscita,
raschiando il fondo della nave, e rami d'albero s'impigliarono nel
carrello di atterraggio. A quanto pare, Skippy le disse di farlo
scendere e salire di nuovo, e ci fu dato il via libera: le porte del
carrello di atterraggio erano chiuse. «Sì. Ora.» Desai premette un
pulsante e sollevò le mani dai comandi. «Pilota automatico attivato.»

Chang e Adams si sporsero per tenersi agli schienali dei nostri sedili.
«Dove stiamo andando?», chiese Chang.

Guardai Desai che parlava con Skippy: «Hacker dice che il pilota
automatico ci sta portando alla base di rifornimento dell'Unef, trecento
chilometri a nord da qui».

«Chi è Hacker?», domandò ancora Chang. «Bishop, abbiamo bisogno di
sapere che cosa sta succedendo. Perché quei ruhar non sono riusciti a
usare i loro fucili?»

«Qui è il colonnello Bishop che ti parla, tenente colonnello Chang.» Ero
incazzato. Non incazzato con lui, incazzato per essere stato beccato a
mentire e non avere una solida via d'uscita. «Hacker è un'unità
informatica dell'Unef», bugia, «mi hanno contattato con un piano per
tirarci fuori di lì», verità, «è parte di un tentativo di contrattacco
alle lucertole», verità di nuovo. Due su tre non è male, stavo andando
alla grande.

«I kristang sono nostri alleati», disse Chang.

«Stronzate. Signore», sparò Desai con rabbia di rimando.

«Chang, se credi ancora che i kristang siano nostri alleati, possiamo
farti scendere da qualche parte.» Avevo bisogno di sapere se Chang
avrebbe rappresentato un problema.

«Non ho creduto che fossero nostri alleati dal primo mese in cui siamo
arrivati qui», spiegò Chang, «ma questo non cambia il fatto che siamo
ufficiali militari e la nostra catena di comando prende ordini dai
kristang. O il fatto che controllano la Terra. E che non abbiamo un modo
per tornare a casa, o contrastare i kristang quando ci arriviamo.»

«Ci stiamo lavorando», dissi con sincerità. «Non so quanto credere a
questo Hacker», vero anche questo, «ciò che so è che, chiunque sia,
siamo arrivati fin qui.» Abbastanza vicino alla verità. «Il passo
successivo sarà caricare i rifornimenti e qualche volontario. Dopo di
che, saliremo a bordo e ruberemo una nave da guerra kristang.»

«Mi sta prendendo in giro, signore.» Adams mi lanciò uno sguardo
incredulo. «Come facciamo a sapere che questo Hacker è\ldots»

«Adams, se ti avessi detto ieri mattina che saremmo evasi di prigione
due volte, avremmo rubato armi ruhar e saremmo decollati con un Dodo, mi
avresti creduto allora? Abbi un po' di fiducia, è tutto quello che
possiamo avere qui.»

Tutti e tre i miei compagni brontolavano a bassa voce, mentre il Dodo
seguiva la sua rotta programmata verso la nostra meta, qualunque fosse.
Ci spaventammo quando una coppia di Avvoltoi passò nella direzione
opposta, con una lieve oscillazione delle ali in segno di saluto. Skippy
doveva aver parlato con loro per radio, perché ci scivolarono proprio
accanto. Questo colpì i miei amici scettici, anche Chang.

«Hai ricevuto ordini da questo Hacker?», chiese Chang sospettoso.
«Perché l'Unef avrebbe assegnato te, o noi, a questa missione? Eravamo
in prigione all'epoca. Non siamo una unità di forze speciali. Non ha
alcun senso.»

«Ci hanno scelto perché eravamo nel posto giusto al momento giusto,
nessun altro motivo. Senti, colonnello Chang, l'Unef è un comando
unificato, ma Hacker è dell'esercito degli Stati Uniti e tu non ne fai
parte.» Chang era un buon ufficiale, non mi piaceva mentirgli, a nessuno
di loro. «Quando arriviamo alla base di rifornimento, puoi scendere da
questa nave, se vuoi. Questa è una missione a titolo volontario.»

«La missione colpirà i kristang?», chiese Desai.

«Li colpirà duro e non potranno prevederlo.» Se Skippy stava dicendo la
verità, tutto quello che avrebbero saputo i kristang era che il wormhole
per la Terra avrebbe smesso di funzionare.

Desai non ebbe esitazioni: «Io ci sto, signore».

«Anch'io.» Adams accettò con rabbia, toccandosi senza volerlo il
bicipite destro, dove avevo visto una brutta cicatrice lasciatale dai
kristang, alla rapida occhiata che le avevo dato mi aveva ricordato
un'ustione elettrica.

Chang rifletté per un momento, poi mi rivolse un sorriso che non seppi
interpretare. «Quando gli dissero che un ufficiale era brillante e
coraggioso, pare che Napoleone abbia replicato: ``Sì, ma è fortunato?''.
Bishop, non so perché, ma in qualche modo hai il talento di essere nel
posto giusto al momento giusto. Mi fiderò di te per questo motivo. Ci
sto anch'io, qualunque sia questa missione.»

Ottimo. Avevo tre volontari, me ne sarebbero bastati altri venti. Venti
sconosciuti, ai quali sarebbe stato chiesto di fidarsi di me dopo che il
nostro Dodo sarebbe caduto davanti ai loro occhi come una manna dal
cielo. Avrei avuto bisogno di rinforzi. «Hacker mi ha ordinato di non
fornire informazioni dettagliate fino a quando la squadra non avrà
lasciato l'orbita, perché, se questa operazione va a rotoli, non
possiamo rischiare che le lucertole sappiano che l'Unef, gli umani erano
coinvolti. Mi fido di voi però al punto di dirvi questo: se avremo
successo, chiuderemo il wormhole che dà alle lucertole, e ai thuranin,
accesso alla Terra. Ciò significa che la Terra tornerà in territorio
ruhar e sarà troppo lontana da un wormhole perché entrambe le parti si
prendano la briga d'inviare navi.» Vidi quanto fossero scioccati i volti
di tutti, sapevo come ci si sentiva. «Questa missione non ha come fine
quello di colpire le lucertole, non è una vendetta», mentre lo dicevo,
guardai negli occhi Desai, «il fine è salvare la Terra.»

«Merda», Adams respirò piano.

«Siamo in grado di farlo? Hai un piano?», chiese Chang.

«C'è un piano, sì. Non saremmo arrivati fin qui senza i mezzi per
violare i sistemi nemici, e uno di questi sistemi è la rete dei
wormhole.»

Il resto del breve volo andò tutto liscio. Mentre il pilota automatico
era in funzione, Skippy cercava di far acquisire a Desai familiarità con
i comandi e noialtri esploravamo la nave alla ricerca di qualsiasi cosa
potesse rivelarsi utile. Sgomberammo tre bidoni che contenevano cibo per
criceti, da buttare dopo l'atterraggio. «Hai capito come far volare
questa cosa, Desai?»

Desai alzò gli occhi verso di me e tenne il pollice in orizzontale, né
su né giù. «I piloti hanno un detto, signore: ``Se riesco a farlo
partire, posso tenerlo in aria''. Nelle intenzioni era una battuta.
Hacker mi ha spiegato i controlli di base, ma è probabile che si sarebbe
schiantato se non fosse inserito il pilota automatico.»

«Puoi riportarci in aria?»

«L'ho fatto una volta», disse lei senza convinzione. Il pilota
automatico fece un atterraggio perfetto, mentre Desai stava a guardare e
mimava l'uso dei comandi.

«Cazzo, i criceti sono già qui, signore», riferì Adams dopo avere
guardato fuori dal finestrino. «Sei criceti armati di fucile, che ci
aspettano.»

Non sorprendeva che i ruhar avessero occupato una grande base di
rifornimento dell'Unef. «Hacker, abbiamo una compagnia ostile qui»,
dissi allo zPhone.

«Ricevuto, Planter, si aspettano una squadra di soldati criceti sul tuo
Dodo. Ci sono solo sei ruhar alla base al momento. Sappiate che le loro
armi sono disattivate, passo.» Skippy si stava divertendo.

Avviammo l'abbassamento della rampa posteriore come distrazione, poi
aprimmo la porta laterale e ci riversammo fuori tutti e quattro, con una
raffica di fasci stordenti. Tutti e sei i criceti caddero in fretta, del
tutto colti di sorpresa. Calciammo via le loro armi, e, li stavamo già
legando, quando un gruppo di umani disarmati ci si avvicinò. Il loro
capo era un maggiore donna dell'esercito americano, Simms, stando al
nome sul suo cartellino. «Che cazzo state facendo? Abbiamo una tregua
con i ruhar!»

Mi alzai e feci il saluto. «Colonnello Joe Bishop. Sì, il tizio di
Barney.» Cavolo, ormai questa storia stava stufando. «Non stiamo
combattendo contro i ruhar, abbiamo requisito una delle loro navi»,
indicai il Dodo, «e abbiamo bisogno di caricare i nostri rifornimenti.
Senza stupide domande dei criceti qui.»

Simms inclinò la testa. Avevo ancora indosso i pantaloni kristang a
strisce nere e gialle, e la giacca della mia divisa era ancora
imbrattata di sangue secco di kristang. Adams e Desai indossavano abiti
altrettanto spaiati e avevano stracci al posto delle scarpe. Il volto di
Chang era ferito, con tagli che i ruhar avevano medicato alla buona. E
io avevo due dita slogate. Forse rotte. «Colonnello Bishop, signore,
l'ultima volta che abbiamo avuto sue notizie era prigioniero dei
kristang.»

«Mi hanno fatto uscire prima per cattiva condotta. È al di fuori della
tua giurisdizione, maggiore. Siamo un'unità di forze speciali, abbiamo
bisogno di provviste e volontari per un assalto ai kristang.»

«Assalto ai kris\ldots{} signore? Non ho ricevuto ordini dal quartier
generale dell'Unef riguardo a unità di forze speciali.»

«Come hai saputo di una tregua?», chiesi. Io non ne avevo saputo nulla.
Ovvio, ero stato in prigione, di nuovo.

Lei tamburellò con le dita l'auricolare dello zPhone. «Annuncio
dall'Unef la scorsa notte. Non siamo riusciti a metterci in contatto per
confermare e questi ruhar sono arrivati qui ieri e hanno preso le nostre
armi.»

«Controlla di nuovo i tuoi messaggi, dovresti aver ricevuto ordine di
assisterci.» Skippy era in ascolto, se era rapido e potente come si
vantava di essere, il maggiore Simms avrebbe presto avuto quel
messaggio.

«Signore?» Alzò gli occhi dallo schermo dello zPhone, sorpresa: «Ho
davvero ricevuto ordini. Come\ldots».

«Le spiegazioni possono aspettare, abbiamo tempi stretti. Quante persone
avete qui?»

«Sessantadue, per lo più corpi di approvvigionamento. Diciotto dei quali
di fanteria, di tutte le nazionalità.»

«Bene. Ho bisogno di una ventina di volontari, o giù di lì, meglio se
con esperienza di combattimento, per una missione delle forze speciali
fuori dai confini del pianeta. E provviste, armi, munizioni, cibo,
medicinali.»

Simms non sembrava ancora del tutto convinta: «Questo è molto insolito.
State per attaccare i kristang, non i ruhar?».

«Maggiore, siamo fanti su un pianeta alieno che cambia gestione a
seconda di chi ha la flotta più grande in orbita al momento, e la cosa
che tu ritieni insolita è una missione delle forze speciali? Gli ordini
che hai ricevuto dal quartier generale dell'Unef contengono i codici di
autenticazione corretti, giusto?» Ovvio che Skippy aveva tutto sotto
controllo, in qualche modo. «Sì? Allora puoi salire a bordo, o toglierti
di mezzo.» Mi rivolsi ai miei compagni: «Colonnello Chang, sergente
scelto Adams, controllate quali rifornimenti hanno qui, ultimate il
carico il prima possibile. Capitano Desai, continua la tua, ehm», è
probabile che non fosse una buona idea usare la parola ``formazione'' di
fronte a persone che speravo si offrissero volontarie, «i tuoi
preparativi di volo». Lei sapeva cosa intendevo. Mi fermai un momento,
mentre Skippy mi parlava piano nell'auricolare. «E questi fucili ruhar»,
ne colpii uno con il piede, «sono di nuovo operativi, perciò li
porteremo con noi. Maggiore Simms, raduna i tuoi.»

Bisogna darle atto che Simms si adattò in fretta. Mi fece un saluto
freddo e trottò via, gridando ordini. Nel giro di cinque minuti, la
maggior parte delle persone alla base stavano gironzolando dietro il
Dodo, io ero in cima alla rampa, in modo che la gente potesse vedermi.
La base non era grande, consisteva per la maggior parte in due lunghi
magazzini, una sorta di prefabbricati di cemento, un paio di fabbricati
annessi, tende per gli umani di stanza lì, piazzole di atterraggio e
un'unica lunga pista. Non era come uno degli enormi magazzini di
rifornimenti dell'Unef, raggruppati attorno alla base dell'ascensore
spaziale, era un centro logistico regionale. A un tratto mi venne in
mente che dovevo fare il discorso della mia vita e non avevo preparato
nulla da dire. «Buongiorno!», esordii a voce alta. «Sono il colonnello
Joe Bishop. Alcuni di voi sanno chi sono. Per quelli di voi che non mi
conoscono, ho catturato un soldato ruhar nella mia città natale nel
Maine, ho abbattuto due Balene ruhar al Vettore di carico e di recente
sono stato prigioniero dei kristang per essermi rifiutato di obbedire
all'ordine di uccidere dei civili criceti.» In questa situazione, un po'
di spavalderia e il fatto di ricordare alla gente perché fossi famoso
tornava a mio vantaggio. «Pensavamo che i kristang fossero i nostri
salvatori, i nostri alleati, quando hanno cacciato i ruhar dopo il loro
assalto sulla Terra. Ora sappiamo che i ruhar hanno colpito la nostra
infrastruttura industriale solo perché i kristang erano alla conquista
del nostro pianeta e i ruhar volevano rovinare loro il bottino di
guerra. Avrete tutti sentito parlare di biscotti della fortuna
provenienti dalla Terra. Beh, non so che notizie di seconda o terza mano
abbiate avuto, perciò ve ne do io di fresche qui. Le lucertole stanno
violentando il nostro pianeta natale. Non so se i ruhar siano potenziali
alleati, o neutrali, o malvagi come i kristang, ma so questo: le
lucertole sono i nostri nemici.» Un forte mormorio si levò dalla folla.
«Quando sono stato promosso, sono andato di sopra per incontrare le
lucertole, una di loro si è ubriacata e ha detto a me e al tenente
colonnello Chang quello che pensano davvero dell'umanità. Le lucertole
ci ritengono dei deboli, dei rammolliti, dei cavernicoli ignoranti buoni
a nulla, se non come soldati semplici e schiavi. L'Unef mi ha fatto
piantare patate, perché le lucertole non vogliono spendere altri soldi
per portare provviste dalla Terra, perché siamo sacrificabili. Le
lucertole avevano già messo in lista me e Chang per il plotone
d'esecuzione. Le donne -- avete sentito cosa pensano i kristang delle
donne -- le nostre donne sono state torturate e stavano per essere
impiccate, quando un'incursione ruhar ci ha liberati dalla prigione.»
Quell'osservazione fece sì che lo sguardo di Simms s'indurisse e la sua
bocca si riducesse a una linea stretta. In quel momento ebbi la certezza
che sarebbe stata entusiasta di qualsiasi cosa avessi voluto proporle.
Adams passò trotterellando tra la folla e salì la rampa fino a me,
facendomi un saluto deciso.

«I rifornimenti sono scarsi, signore, ma adeguati.»

«Bene. Vedo che hai trovato degli stivali.»

Lei si lanciò un'occhiata ai piedi con un ampio sorriso. «E pantaloni,
signore», aggiunse con uno sguardo significativo ai miei larghi
pantaloni kristang.

«Ascoltate, gente!», mi rivolsi di nuovo alla folla. «Questo non è un
altro assalto dei ruhar, stavolta i criceti sono qui per restare.
Sappiamo per certo dall'intelligence che i ruhar e i loro alleati hanno
sconfitto un gruppo combinato thuranin-kristang e che i thuranin si
stanno ritirando da quest'area; non stanno più sostenendo lo sforzo dei
kristang per mantenere questo pianeta. Ciò significa che la missione per
cui l'Unef è venuta qui è finita e che non abbiamo modo di ottenere
rifornimenti dalla Terra, o di tornare sulla Terra. Siamo tagliati
fuori. La nuova missione dell'Unef è la sopravvivenza: o piantiamo
colture e ci occupiamo dei raccolti o moriamo di fame. Gli esseri umani
su questo pianeta sono agricoltori ora, non soldati.» Trattandosi di una
base logistica, è probabile che le persone sul posto avessero constatato
gli effetti del calo delle nostre spedizioni di rifornimenti prima che
qualcuno sul campo se ne accorgesse.

«L'Unef sta organizzando una missione rapida delle forze speciali per
colpire i kristang, abbiamo bisogno di volontari. Alcuni di voi hanno
visto che i criceti non sono stati in grado di usare le armi quando
siamo atterrati. È l'Unef, hanno un modo per hackerare i sistemi ruhar e
kristang.» Era abbastanza vero, io facevo parte dell'Unef e il mio modo
di hackerare era chiedere a Skippy di farlo. Dovevo stare attento a
quello che dicevo sulla missione, i ruhar avrebbero interrogato a fondo
le persone che ci saremmo lasciati alle spalle. «Abbiamo requisito un
Dodo e sferrato un attacco di spoofing ai sistemi di controllo del
traffico aereo ruhar, così non ci vedranno. L'Unef non sa quanto durerà
questa finestra di opportunità, perciò abbiamo un programma serrato.
Ecco cosa posso dirvi dell'operazione: se funziona, andremo in orbita,
oppure oltre, per colpire duro i kristang. Le persone che verranno con
noi riceveranno informazioni dettagliate una volta che avremo lasciato
l'atmosfera. Questa è un'occasione per fare la vera differenza in questa
guerra.» Se mi fossero spuntate le ali e avessi preso il volo, la gente
sarebbe stata meno sorpresa. «La missione di combattimento qui», indicai
per terra, «è finita, se volete un'opportunità di contrattacco alle
lucertole, è con noi.»

Feci una pausa per controllare i volti della folla. Era tutto troppo,
troppo veloce. Non molto tempo prima, eravamo tutti beatamente soli
nell'universo, poi la Terra era stata attaccata, si era formata l'Unef e
noi eravamo stati portati in fretta e furia su un altro pianeta. Fino a
qualche mese prima, eravamo concentrati sull'avere un buon rendimento in
una missione difficile per i nostri alleati, poi erano arrivati i
biscotti della fortuna e avevamo scoperto che i kristang non erano amici
dell'umanità. E, dopo il grande assalto dei ruhar quando ero al Vettore,
avevamo saputo che i kristang non erano in grado di garantire la nostra
sicurezza su Paradiso. Soltanto la mattina precedente, le lucertole
erano ancora saldamente a capo del pianeta, ora stavo dicendo a quegli
uomini e a quelle donne che quella condizione era finita, per sempre. E
che l'Unef sarebbe rimasta bloccata lì per molto tempo. E che, in
qualche modo, come per un miracolo di sospetta tempestività, avevo un
piano per consentirci di lasciare il pianeta e contrattaccare. Se fossi
stato tra la folla ad ascoltare un buffone che mi raccontava tutto
questo, per me non si sarebbe trattato d'altro che di un mucchio di
stronzate. Vedevo che le persone si agitavano, bisbigliandosi cose a
vicenda, cercando di decidere il da farsi.

Qualcuno chiamò a voce alta in cinese, o almeno io diedi per scontato
che fosse cinese; Chang era tornato dal magazzino. Tre soldati in
uniforme dell'Esercito popolare di liberazione si girarono quando Chang
parlò. In quel momento compresi che forse non avevano capito una sola
parola di quello che avevo detto. Chang si fece strada tra la folla per
raggiungerli, parlarono in cinese per breve tempo, poi tutti e quattro
raggiunsero a piedi la base della rampa. «Altri tre volontari,
colonnello Bishop.»

Questo mi mise in una situazione scomoda, gli feci cenno di salire la
rampa. «Colonnello Chang, questa è una missione volontaria, non voglio
che ordini alla gente di partire. Cosa hai detto loro?»

Chang sbatté le palpebre, sorpreso: «Ho detto che questa missione sarà
l'unica opportunità che avranno per fare il loro dovere e servire il
popolo cinese. Si sono offerti volontari».

Ah, che diavolo. Aveva detto la verità. La missione era importante e
avevamo bisogno di soldati. «Non male.»

«Merda», disse Adams a voce alta. «Signore, dobbiamo andare in una base
che ha veri soldati e Marine, non questi passacarte.»

Forse il mio discorso era stato efficace e la gente si sentiva
intrappolata su Paradiso e voleva un modo per fare qualcosa, e la folla
aveva solo bisogno di una spinta in più. Forse vedere tutti e tre i
cinesi unirsi a noi li motivò. Forse tutto ciò di cui avevano sempre
avuto bisogno per indurli a muoversi era che un Marine li facesse
vergognare.

«Oh, cazzo no. Non mi farò mettere i piedi in testa da un Marine.» Il
maggiore Simms si fece avanti. Dopo di lei, ci fu un'ondata di consensi.
Dei diciotto fanti, diciassette si offrirono volontari; uno aveva un
amore su Paradiso e non voleva lasciarlo lì a un ignoto futuro. Presi
tutti i diciassette fanti: la nostra unità di forze speciali ad hoc era
forte di ventiquattro persone. Ventiquattro persone e una lucida lattina
di birra parlante.